\section{Introduction}
\label{sec:intro}

Salient object detection (SOD) seeks to segment the single most visually
distinctive object in a scene, producing a binary, class-agnostic mask. Inspired
by the human visual system's ability to instantly attend to the most informative
region of a scene, SOD is a long-standing pre-processing step for a wide range of
downstream applications, including image retrieval, object-aware image and video
editing, visual tracking, image captioning, and content-aware
compression~\cite{borji2015salient,wang2021salientsurvey}, where a compact
foreground estimate is far more useful than a dense semantic labelling. Because the
task is defined by perceptual distinctiveness rather than a fixed set of
categories, a good SOD model must reason jointly about local contrast, global
scene context, and a generic, category-independent notion of ``objectness''.

Modern SOD networks are almost universally fully-supervised encoder--decoder
architectures trained from scratch or from ImageNet-pretrained
backbones~\cite{basnet2019,f3net2020,ldf2020,vst2021,edn2022,menet2023}. Early
CNN-based methods introduced multi-scale aggregation and contextual attention to
capture both fine boundaries and global structure~\cite{picanet2018}, and a large
body of subsequent work focused on boundary quality through hybrid region--edge
losses, feedback refinement, and body--edge
decoupling~\cite{basnet2019,f3net2020,ldf2020}. More recently, transformer
backbones have pushed accuracy further by modelling long-range dependencies among
image tokens~\cite{vst2021,menet2023}, while a parallel line of lightweight
designs targets real-time inference on resource-constrained
devices~\cite{csnet2020,edn2022,sod_admnet,sod_bdfgnet,sod_stnet}. Despite their
architectural diversity, these methods share a common limitation: their inductive
biases and their notion of objectness are learned entirely from the relatively
small saliency training sets, and they carry no external semantic knowledge about
the objects they are asked to segment.

In parallel, vision foundation models have transformed segmentation. The
Segment Anything Model (SAM)~\cite{kirillov2023segment} provides a strong,
promptable class-agnostic segmenter trained on billion-scale masks, while
CLIP~\cite{radford2021learning} supplies open-vocabulary vision--language
alignment learned from web-scale image--text pairs. Recent open-vocabulary
segmentation networks such as CAT-Seg~\cite{cho2024cat} and, most relevantly,
ESC-Net~\cite{lee2024escnet} couple the two by building a vision--language
\emph{correlation map} between dense CLIP features and text-class embeddings and
using it to drive SAM's prompt decoder. These systems inherit the semantic
knowledge of CLIP and the boundary precision of SAM without training either from
scratch, and they achieve strong open-vocabulary segmentation with only a small
number of trainable fusion parameters. This raises a natural question:
\emph{can the SAM--CLIP correlation paradigm be adapted to give SOD the same
foundation-model benefits?}

The adaptation is not trivial. Open-vocabulary segmenters are designed for
multi-class recognition: each class contributes its own text prompt, and the
correlation map is naturally aligned to a specific semantic label. SOD is
class-agnostic and binary: there is no fixed class to prompt. A naive port that
uses a single generic prompt such as ``a salient object'' aligns poorly with the
enormous appearance diversity of salient objects and collapses for small or
low-contrast targets, producing empty or diffuse correlation maps that in turn
yield degenerate SAM prompts. The pseudo-prompt generator that converts the
correlation into SAM inputs is especially brittle: after a spatial softmax the map
is nearly uniform, so a fixed threshold routinely produces an empty prompt set and
starves the SAM decoder of spatial guidance. A coarse, multi-class-style decoder
compounds the problem by discarding the fine detail that SOD demands.

%% OLD (submitted): "We investigate what it takes to make this adaptation work, and
%% reach a perhaps surprising conclusion: the dominant factor is ... the gradient
%% path ... Stacking multiple ESC blocks progressively overwrites the initial prompt
%% correlation ... A prompt-residual connection ... restores this signal, and only
%% then does the prompt mechanism pay off. ... a plain U-Net decoder outperforms the
%% complex ASPP/FPN decoder and the auxiliary point-classification loss ..."
\rev{We investigate what it takes to make this adaptation work. With a CLIP
ViT-L/14 backbone, a prompt bank, and a pseudo-prompt generator that never returns
an empty prompt set, a simple configuration (U-Net decoder, binary cross-entropy)
reaches strong results on all benchmarks, and adding a correlation-aware ASPP/FPN
decoder or an auxiliary point-classification loss does not improve S-measure on
the primary benchmark. We also examine how the prompt-derived correlation reaches
the loss. In the binary setting the text cross-attention inside the fusion block
has a single key, so its output does not depend on its inputs; the prompt branch
then receives gradient only through a residual that adds the initial correlation
back before decoding. In matched
experiments this residual does not improve accuracy, and repeated runs show that
several component differences are within seed-to-seed variance. By contrast,
fine-tuning the CLIP visual encoder matters a great deal: freezing it lowers
S-measure by $0.024$ on DUTS-TE and $0.044$ on DUT-OMRON. }

%% OLD (submitted) contributions: (1) prompt-residual correlation architecture ...
%% "without this residual, the VLF blocks completely overwrite the prompt-derived
%% spatial prior"; (2) learned prompt-attention aggregation ... "improving S-measure
%% on the challenging SIP benchmark"; (3) adaptive-fallback PPG ... "the single most
%% impactful component in our ablation"; plus "a comprehensive analysis ...".
\rev{\noindent\textbf{Contributions.}
\begin{itemize}
  \item An adaptation of the SAM--CLIP correlation paradigm to binary SOD with a
    CLIP ViT-L/14 backbone, a fixed bank of $K{=}10$ salient-object prompts, and an
    \emph{adaptive-fallback pseudo-prompt generator} that guarantees non-empty
    spatial guidance for SAM. Among the compared methods it obtains the best
    S-measure on four of five RGB benchmarks and, with RGB input only, outperforms
    the listed RGB-D methods on SIP.
  \item A \emph{diagnostic analysis of the prompt gradient path}: we show that the
    single-key text cross-attention of the fusion block makes its output
    independent of the correlation, measure the gradient reaching the prompt branch
    with and without a prompt-residual connection, and report a matched comparison
    showing that the residual restores the gradient but does not improve accuracy.
  \item An analysis of prompt aggregation and bank design over repeated runs:
    with the full CLIP text encoder, max-pooling and learned per-pixel attention,
    bank sizes $K\in\{1,5,10\}$, and three prompt-wording strategies differ by at
    most $0.006$ S-measure on the RGB benchmarks ($0.010$ on SIP), while fine-tuning
    the CLIP visual encoder is essential (freezing it costs up to $0.044$
    S-measure).
\end{itemize}}
